% !TeX program = xelatex
\documentclass[UTF8,12pt,a4paper]{ctexart}
\usepackage[margin=2.6cm]{geometry}
\usepackage{booktabs,tikz,tabularx,array,enumitem,fancyhdr,xcolor}
\usepackage[colorlinks=true,linkcolor=blue!55!black]{hyperref}
\newcolumntype{Y}{>{\raggedright\arraybackslash}X}
\newcommand{\ubrk}{\textunderscore\allowbreak}
\newcommand{\fname}[1]{{\ttfamily\let\_\ubrk #1}}
\newcommand{\key}[1]{\textbf{\textcolor{blue!45!black}{#1}}}
\newcommand{\why}[1]{\par\smallskip\noindent\textcolor{orange!45!black}{\textbf{为什么低：}}#1\par\smallskip}
\setlist{nosep,leftmargin=1.8em}
\renewcommand{\arraystretch}{1.3}
\linespread{1.12}
\pagestyle{fancy}\fancyhf{}
\fancyhead[L]{\small 重庆青少年心理健康项目}
\fancyhead[R]{\small 阶段汇报}
\fancyfoot[C]{\thepage}
\setlength{\headheight}{15pt}
\ctexset{section={format=\Large\bfseries},subsection={format=\large\bfseries}}

\title{\textbf{重庆青少年心理健康项目}\\[0.5em]阶段汇报\\[0.4em]\large 五种客观测量的目的、结论与原因分析}
\author{}
\date{报告区间：2026 年 9 月 7 日 -- 9 月 10 日}

\begin{document}
\maketitle
\thispagestyle{fancy}

\section*{一页纸结论}
\addcontentsline{toc}{section}{一页纸结论}

这四天里，我们把手上\key{全部五种客观测量}——脑电、近红外、面部视频、行为按键、眼动——都完整地试了一遍，并且第一次对脑电动用了深度学习。

\key{主要结论：五种测量都没能超过“问三个问题”的效果。}那三个问题是年龄、性别、年级。只用这三项建立的模型，判断水平已达 0.67（后文解释这个数字怎么读）。

但这句话如果只看一半，会得出错误的印象。有四点请导师特别注意：

\begin{itemize}
\item \key{测量本身是好的。}五种数据里都验证到了教科书上该出现的生理现象——脑电的 P300 波、大脑血流的正常反应、孩子看到笑脸会去看嘴巴、做错题之后反应变慢。数据、仪器、处理流程都没问题。

\item \key{“低”是有具体原因的，而且原因各不相同。}有些是采集环节真的丢了数据（比如一台设备的按键从头到尾没记上），有些是实验设计本身不够（比如每种情绪只放了 12 张照片，导致指标根本测不准），有些是任务测的东西和情绪无关。第 \ref{sec:why} 节逐条列出。这些是可以改进的。

\item \key{脑电深度模型最接近成功。}在放宽判断标准后，脑电有 35 个结果存活下来，是五种测量里最多的。第 \ref{sec:eeg} 节完整交代这批结果——包括它们为什么最终没被采信。

\item \key{本阶段最重要的产出是一套判断规则。}四天里出现过三次“看起来成功了”的结果，全部被这套规则拦下。规则是在真实错误里一条条长出来的。
\end{itemize}

\vspace{0.5em}
\begin{center}
\begin{tabularx}{\textwidth}{@{}lY@{}}
\toprule
\textbf{问题} & \textbf{回答} \\
\midrule
客观测量能诊断吗？ & 在本队列、按目前方法，不能超过年龄性别年级。 \\
是数据本身没信息吗？ & 不是。同一套方法猜年龄能到 0.82、猜性别 0.77。 \\
那为什么低？ & 采集缺陷、试次不足、设备与学校混淆、任务与情绪无关。见第 \ref{sec:why} 节。 \\
有接近成功的吗？ & 有。脑电深度模型，见第 \ref{sec:eeg} 节。 \\
下一步做什么？ & 见第 \ref{sec:next} 节，四个方向。 \\
\bottomrule
\end{tabularx}
\end{center}

\newpage
\section{这个项目在问什么}

\subsection{任务}

我们收集了三千多名中小学生的多种客观测量数据，想回答：\key{能不能不依赖问卷和访谈，只靠仪器测到的信号，判断一个孩子是否存在心理健康问题（有／无）。}

\subsection{为什么“年龄性别年级”是一道必须跨过的门槛}

这是理解整份报告的关键。

在我们这批孩子里，心理问题的出现率本身就和年龄、性别、年级相关——年级越高比例越高。所以哪怕完全不看仪器数据，只知道“这是一个初三的女生”，就已经能做出相当程度的猜测。

我们实测了这个水平：\key{只用年龄、性别、年级三项，达到 0.67。}

因此任何一种客观测量要证明自己有用，\key{光比瞎猜强是不够的}，必须证明：在已经知道年龄性别年级的前提下，再加上它，判断会更准。这个“额外提升”是全文的核心概念。

\subsection{这些数字怎么读}

\begin{center}
\begin{tabularx}{0.92\textwidth}{@{}lY@{}}
\toprule
\textbf{数值} & \textbf{含义} \\
\midrule
0.50 & 完全等于抛硬币，没有任何判断能力 \\
0.60 & 有一点点能力，但很弱 \\
0.67 & 年龄性别年级三项达到的水平 \\
0.80 & 相当不错 \\
1.00 & 完美，永不出错 \\
\bottomrule
\end{tabularx}
\end{center}

\noindent 更直观的说法：\key{0.67 的意思是，随便挑一个有问题的孩子和一个没问题的孩子，模型有 67\% 的把握把前者排在前面。}

\section{五种测量的目的与结论}

\subsection{总览}

下表每一行都是\key{在同一批孩子身上、用同一个模型}算出来的三个数：这个模态自己有多强、只用年龄性别年级有多强、以及\key{把两者合起来}有多强。

最后一列是本报告最关键的一个量——\key{把模态加进去之后，判断到底改善了多少}，方括号里是它的 95\% 置信区间。

\vspace{0.4em}
\begin{center}
{\small\setlength{\tabcolsep}{4pt}
\begin{tabular}{@{}llcccl@{}}
\toprule
\textbf{测量} & \textbf{人数} & \textbf{模态单独} & \textbf{只用人口学} & \textbf{人口学+模态} & \textbf{提升与 95\% 区间} \\
\midrule
眼动 & 305 & 0.578 & 0.653 & 0.674 & +0.021\quad[-0.048, +0.089] \\
行为按键 & 1412 & 0.558 & 0.590 & 0.608 & +0.018\quad[-0.005, +0.040] \\
近红外 & 524 & 0.547 & 0.579 & 0.592 & +0.013\quad[-0.041, +0.067] \\
脑电（深度） & 1820 & 0.534 & 0.602 & 0.604 & +0.002\quad[-0.006, +0.008] \\
面部视频 & 3381 & 0.617 & 0.669 & 0.666 & -0.002\quad[-0.017, +0.011] \\
脑电（特征） & 1820 & 0.527 & 0.604 & 0.571 & -0.033\quad[-0.059, -0.006] \\
\bottomrule
\end{tabular}}
\end{center}

\vspace{0.6em}
\begin{center}
\begin{tikzpicture}[x=1cm,y=1cm]
\def\xs{34}
\def\dy{1.15}
\foreach \i/\name/\sig/\dem/\comb in {
  0/眼动/0.578/0.653/0.674,
  1/行为按键/0.558/0.590/0.608,
  2/面部视频/0.617/0.669/0.666,
  3/近红外/0.547/0.579/0.592,
  4/脑电（深度）/0.534/0.602/0.604,
  5/脑电（特征）/0.527/0.604/0.571}{
  \node[anchor=east,font=\small] at (-0.1,\i*\dy) {\name};
  \pgfmathsetmacro\ls{(\sig-0.50)*\xs}
  \pgfmathsetmacro\ld{(\dem-0.50)*\xs}
  \pgfmathsetmacro\lc{(\comb-0.50)*\xs}
  \fill[blue!30]   (0,\i*\dy+0.30) rectangle (\ls,\i*\dy+0.48);
  \fill[orange!55] (0,\i*\dy+0.09) rectangle (\ld,\i*\dy+0.27);
  \fill[green!45!black!60] (0,\i*\dy-0.12) rectangle (\lc,\i*\dy+0.06);
  \node[anchor=west,font=\scriptsize] at (\ls+0.08,\i*\dy+0.39) {\sig};
  \node[anchor=west,font=\scriptsize] at (\ld+0.08,\i*\dy+0.18) {\dem};
  \node[anchor=west,font=\scriptsize] at (\lc+0.08,\i*\dy-0.03) {\comb};
}
\draw[thick,gray] (0,-0.45) -- (0,6.4);
\node[anchor=north,font=\scriptsize,gray] at (0,-0.5) {0.50（抛硬币）};
\fill[blue!30]   (7.3,6.30) rectangle (7.75,6.46);
\node[anchor=west,font=\scriptsize] at (7.8,6.38) {模态单独};
\fill[orange!55] (7.3,6.02) rectangle (7.75,6.18);
\node[anchor=west,font=\scriptsize] at (7.8,6.10) {只用人口学};
\fill[green!45!black!60] (7.3,5.74) rectangle (7.75,5.90);
\node[anchor=west,font=\scriptsize] at (7.8,5.82) {人口学+模态};
\end{tikzpicture}
\end{center}

\vspace{0.3em}
\noindent 图里有两件事值得看。

\key{第一，每一根蓝条都是三根里最短的。}任何一个模态单独拿出来，都不如年龄性别年级。但\key{短的原因各不相同}，这正是第 \ref{sec:why} 节要讲的。

\key{第二，绿条和橙条几乎一样长。}把模态加进去之后，成绩基本没动——这就是全部结论所在。

\subsection{为什么“绿条比橙条长一点”还不够}

看表格最后一列：\key{六组里有四组的提升是正的}（眼动 +0.021、行为 +0.018、近红外 +0.013、深度脑电 +0.002）。这很容易让人得出“那这些模态还是有用的”这一结论。

但这个结论下不得，原因在方括号里。

方括号里是\key{置信区间}，表示这个差值真实可能落在什么范围内。判断标准只有一条：\key{区间必须整个在零的一侧。}

看眼动那一行：+0.021 听上去不错，但区间是 $[-0.048, +0.089]$——\key{从“明显有害”一直横跨到“明显有用”}。305 个孩子的数据量，分辨不出这 0.021 是真实效应还是随机波动。

打个比方：抛 10 次硬币出现 6 次正面，数字确实大于 5，但没人会因此断定硬币有偏。要下这个判断，得抛几百次。\key{上表六行，没有任何一行的区间排除了零}（唯一排除零的是最后一行，而它在零的负侧）。

这正是第 \ref{sec:false} 节要讲的那三次假阳性的成因——当时就是只看了点估计，没看区间。

\subsubsection*{最后一行的负号不代表“有害”}

脑电传统特征那一行是显著变差的，这一点需要单独解释，否则容易被误读。

原因是\key{“稀释”}：模型要同时消化 3 列人口学变量和 321 列脑电特征，当后者不含有用信息时，前者的作用会被摊薄。这是分析方法的性质，\key{不是脑电把结果搞砸了}。

证据是：Goal 3 改用了一种不会稀释的组合方式（先各自算出预测分数再合并），同样的比较做了 56 次，\key{结果是 0 次显著正、0 次显著负}，全部落在 -0.008 到 +0.005 之间。稀释一消失，负号也就消失了。

\key{所以这里的负号意思是“没有信号”，不是“起了反作用”。}

\subsection{脑电}

\key{目的：}孩子戴电极帽完成“听声音数目标音”的任务（Oddball），检验大脑对刺激的电反应能否区分有无心理问题。

\key{做了什么：}此前的分析有一个严重错误——程序只保留了目标音的反应，把作为对照的标准音全丢了，等于只有一半数据。这次从原始记录重新恢复了完整的两种条件，看到了非常漂亮的 P300 波，在 88\% 的孩子身上都存在。随后第一次对它使用深度学习。

\key{结论：}没有额外提升。在已知年龄性别年级基础上加入脑电深度模型，提升 \key{+0.0016}，几乎为零。但脑电是五种测量里最接近成功的一种，详见第 \ref{sec:eeg} 节。

\key{一个关键对照：}用完全相同的程序、孩子、数据，只换要猜的目标：

\begin{center}
\begin{tabular}{@{}lc@{}}
\toprule
\textbf{让同一套程序去猜} & \textbf{成绩} \\
\midrule
孩子的年龄 & 0.82 \\
孩子的性别 & 0.77 \\
这一声是目标音还是标准音 & 0.82 \\
\textbf{孩子有没有心理问题} & \textbf{0.55} \\
\bottomrule
\end{tabular}
\end{center}

\noindent 程序显然有学习能力。学不到的是“心理状态”。

\subsection{近红外}

\key{目的：}近红外测大脑皮层血流变化，反映脑区在完成任务时的活跃程度。孩子做了五个任务（静息、词语流畅、记忆、听声音、赌博游戏）。

\key{做了什么：}此前卡在技术错误上——读取程序只看文件头，没真正读进数据，导致所有任务反应被判为“不可用”。这次重写了读取程序，五个任务在两台设备上全部解锁，并验证到了正常的血流反应（做词语流畅任务时左脑更活跃，符合语言功能的常识）。

\key{结论：}没有额外提升，而且是五种里最彻底的——即便把判断标准大幅放宽，依然一条都没通过。

\subsection{面部视频}

\key{目的：}孩子观看情绪影片时被摄像头录下，检验面部表情能否反映心理状态。

\key{做了什么：}此前是在 11 分钟录像里均匀取 16 帧，但这 11 分钟混杂了看影片、朗读、答题等完全不同的环节，取出来的画面大部分只反映“孩子长什么样、房间什么样”。这次按实验记录切成了对应三段情绪影片的片段，并计算\key{同一个孩子看负性影片与正性影片时的差异}。

\key{结论：}没有额外提升。但这里有\key{本项目唯一真正站得住的正面发现}。

我们设了一个“背景对照”：把画面里的人脸挖掉，只用背景预测。如果模型靠背景就能预测，说明它学的是学校和房间，不是孩子。结果：\key{面部的“负性减正性”差异显著赢过了自己的背景对照，最强 +0.115。}

这个比较有力，是因为同一个孩子的长相、房间、摄像头、学校在这个减法里\key{自动抵消了}。所以这部分信号确实来自表情本身。\key{只是它仍赢不过年龄性别年级。}

\subsection{行为按键}

\key{目的：}孩子做任务时的每次按键（对错、快慢）都被记录，此前从没人读过。检验反应速度、注意力维持、做错后的调整能否反映心理状态。

\key{做了什么：}第一次系统读取这些日志，提取工作记忆表现、反应时间及波动、错误后减速、赌博任务的输赢反应等。

\key{结论：}没有额外提升，但\key{行为指标本身非常健康}——两台设备上的孩子都表现出教科书式规律：做错一题后下一题会慢 100 毫秒左右（76\% 的孩子如此），第二组比第一组更快。这些规律在两批完全不同的孩子身上独立重现。

\subsection{眼动}

\key{目的：}孩子做了三个任务：自由观看情绪面孔、按指令看向／看开目标、用眼睛追踪移动的点。心理学文献普遍认为抑郁的人看负性面孔更久，我们要检验这一点。

\key{做了什么：}先系统查阅近年文献、写下方法和预期，然后才开始分析。过程中修正了四个技术问题，其中一个很关键：\key{三台设备记录的视线位置都系统性偏低}，不校正的话“人看脸时先看眼睛”这条铁律会在两台设备上反号。

\key{结论：}没有额外提升。三台不同设备、三种采样率、三批几乎不重合的孩子，结论完全一致。

\subsection{一个来自医院的独立验证}

眼动分析全部完成之后，医院才补送来眼动范式的原始说明文档。这是难得的检验机会——我们此前所有参数都是从记录数据里反推的，没有任何文档参照。

\key{结果：文档逐项确认了我们反推的参数。}包括面孔呈现 4 秒、每种情绪 12 张、扫视任务 8 个试次、追踪任务快慢两档恰好是二倍频关系。其中两个参数是我们分析过程中因为“结果不合常理”而主动修正过的，文档证明我们改对了。

\newpage
\section{为什么都低：逐条原因}\label{sec:why}

这是本报告最有实际价值的一节。“没有超过人口学”是一个结果，但\key{原因不止一个，而且大部分是可以改进的}。我们把查到的原因分成四类。

\subsection{第一类：采集环节真的丢了数据}

这些是硬性的数据损失，不是分析方法的问题。

\begin{center}
\begin{tabularx}{\textwidth}{@{}>{\hsize=0.85\hsize\raggedright\arraybackslash}X>{\hsize=1.15\hsize\raggedright\arraybackslash}X@{}}
\toprule
\textbf{问题} & \textbf{后果} \\
\midrule
必可明设备的“听声音”任务\key{从头到尾没有采集到按键} & 640 名孩子里 549 名的按键记录完全是空的，厂商自己的样例记录也是空的。整个单元只能作废。 \\
益瑞德设备的赌博任务\key{反应窗口设成 1 秒，而门在屏幕上停留 4.5 秒} & 只有 38\% 的选择被记录下来（另一台设备是 98\%）。可用的“赢了继续／输了改选”配对从 24 组降到 6 组。 \\
50 名孩子运行的记忆任务版本\key{里根本没有重复出现的图片} & 无法计算“认出重复”这一核心指标，只能排除。 \\
必可明设备的词语流畅任务\key{没有记录任何时间标记} & 无法确定任务何时开始，该任务只能作参考。 \\
\bottomrule
\end{tabularx}
\end{center}

\noindent\key{建议：}前两项应尽快反馈给设备方与实验设计方。

\subsection{第二类：实验设计本身的试次不够}

这一类最值得重视，因为它意味着\key{有些指标从原理上就测不准，再多的分析也救不回来}。

\key{最典型的例子是眼动的情绪面孔任务。}这个任务是整个眼动分析的理论核心——文献认为抑郁的人看负性面孔更久。但这个范式\key{每种情绪只放了 12 张照片}。

我们直接检验了这个指标测得准不准：把每个孩子的试次分成单双数两半，看两半算出来的结果是否一致。

\begin{center}
\begin{tabular}{@{}lcc@{}}
\toprule
\textbf{指标类型} & \textbf{一致性} & \textbf{三台设备都可靠的指标数} \\
\midrule
看某类面孔的绝对时长 & 0.70（良好） & 53 个 \\
\textbf{负性减中性的差值} & \textbf{-0.05（等于没有）} & \textbf{0 个} \\
\bottomrule
\end{tabular}
\end{center}

\noindent 原因是纯粹的算术：两个各由 12 张照片估出来的数相减，误差全部留下，而真实差异反而抵消掉了。三台设备上独立重现了同一现象。

\key{这一点必须说清楚：我们的阴性结果说明的是“这个实验设计测不出情绪注意偏向”，而不是“情绪注意偏向不存在”。}这是关于仪器的结论，不是关于疾病的结论。

同类问题还有两处：

\begin{itemize}
\item \key{扫视任务每组只有 8 个试次。}错误率的最小分辨单位因此是 1/8 = 0.125，而文献中常用的同类任务有 40 到 100 个试次。
\item \key{面部的“负性减正性”差异，幅度只有原始画面信息的 17\%。}减法去掉了长相和房间（这是好事），但剩下的真实表情信号本身就很微弱。
\end{itemize}

\subsection{第三类：设备与学校分不开}

\key{近红外的两台设备，几乎是在两批完全不同的孩子身上使用的——1527 人对 1022 人，重合的只有 1 人。}眼动的三台设备也是各自负责不同的学校。

这带来一个无法在事后修复的问题：\key{“设备差异”和“学校差异”完全混在一起}，无法分辨一个信号到底来自仪器还是来自学校。因此两台设备的数据不能合并分析，样本量被硬生生切成两半。

我们也确认了学校本身就是一个很强的预测因素——\key{光凭孩子来自哪所学校，预测水平就能到 0.67}，和年龄性别年级相当。这说明不同学校的孩子患病比例差别很大，任何分析都必须先把这个因素排除掉，否则会把“学校差异”误认成“生理差异”。

\key{建议：}后续如需扩展数据收集，应让每台设备在每所学校都采集一部分，而不是一台设备包一所学校。

\subsection{第四类：任务测的东西可能和情绪无关}

\begin{itemize}
\item \key{脑电用的 Oddball 任务，测的是注意力和对新异刺激的反应}，本身不是情绪任务。它能非常可靠地反映“这个人有没有在专心听”，但未必反映心情。
\item \key{眼动的追踪任务，其原始文献针对的是精神分裂症}，不是抑郁。这套范式是从另一个临床问题上借用过来的。
\item \key{诊断标签本身可能过于笼统。}目前把抑郁、焦虑、高风险合并成了“有问题”这一类。不同问题的生理表现很可能不同，合并之后互相抵消。
\end{itemize}

\subsection{第五类：样本量与方法层面的限制}

\begin{itemize}
\item 部分队列偏小：近红外的听声音任务只有 524 人，眼动各单元 247 到 336 人，行为学的三任务交集只有 342 人。样本越小，结果越不稳定（第 \ref{sec:false} 节会说明这带来的具体危害）。
\item 深度模型的成绩在不同随机初始值下波动可达 0.05，而我们要检测的效应只有 0.005 量级。\key{这意味着非常微小的效应在当前条件下无法被可靠检出。}
\item 用于最终验证的那批孩子（900 人）在项目早期已被使用过，因此严格意义上的最终验证需要一批全新的孩子。
\end{itemize}

\newpage
\section{最接近成功的一次：脑电深度模型}\label{sec:eeg}

前面说“五种都没超过”，但并非五种都一样地没超过。\key{脑电明显比其他四种更接近。}这一节完整交代这批结果，包括我们最终没有采信它的理由。

\subsection{放宽标准后，脑电存活最多}

我们的判断规则一共九条。为了确认这些规则不是“过于严格以至于埋没了真信号”，我们做了一次压力测试：\key{去掉其中六条，重新检查此前全部 878 次比较。}

\begin{center}
\begin{tabular}{@{}lccc@{}}
\toprule
\textbf{测量} & \textbf{比较次数} & \textbf{放宽后存活} & \textbf{再加一条折数要求} \\
\midrule
\textbf{脑电} & 302 & \textbf{35} & \textbf{20} \\
面部视频 & 66 & 12 & 10 \\
眼动 & 288 & 0 & 0 \\
行为按键 & 84 & 0 & 0 \\
近红外 & 138 & 0 & 0 \\
\bottomrule
\end{tabular}
\end{center}

\noindent \key{脑电存活 35 个，是五种测量里最多的，也是唯一一个在深度模型上还有东西剩下的。}眼动、行为、近红外全部归零。面部存活的 12 个来自已被取代的旧分析。

这 35 个结果里，最强的一批提升在 \key{+0.028 到 +0.031}，统计区间都不含零。如果只看这一页，结论会是“脑电深度模型有效”。

\subsection{但它们集中在一个随机种子上}

深度学习每次训练都要从一个随机初始值开始。我们跑了三个不同的初始值（称为“种子”）。这 35 个结果的分布是：

\begin{center}
\begin{tabular}{@{}lc@{}}
\toprule
\textbf{位置} & \textbf{存活数} \\
\midrule
\textbf{种子 2，分组验证} & \textbf{28} \\
种子 0，标准验证 & 6 \\
种子 0，分组验证 & 1 \\
种子 1 & 0 \\
\bottomrule
\end{tabular}
\end{center}

\noindent 28 个挤在同一格里。把那一格拆开看，情况就清楚了：

\begin{center}
\begin{tabular}{@{}lcccc@{}}
\toprule
\textbf{种子} & \textbf{人口学} & \textbf{脑电深度} & \textbf{人口学+脑电} & \textbf{提升} \\
\midrule
0 & 0.5814 & 0.5499 & 0.5873 & +0.006 \\
1 & 0.6013 & 0.5597 & 0.5812 & -0.020 \\
\textbf{2} & \textbf{0.5574} & \textbf{0.5234} & 0.5865 & \textbf{+0.029} \\
\midrule
三次平均 & 0.5909 & 0.5475 & 0.5878 & -0.003 \\
\bottomrule
\end{tabular}
\end{center}

\noindent 请注意第四列：\key{“人口学+脑电”这一列三次几乎不动}（0.5873、0.5812、0.5865，相差不到 0.006）。真正在跳的是第二列\key{“人口学”自己}，从 0.6013 掉到 0.5574，跳了 0.044。

而且在种子 2 那一次，\key{脑电自己的成绩反而是三次里最差的}（0.5234）。

\subsection{结论}

所以种子 2 的那个“+0.029 提升”，\key{不是脑电变强了，是被比较的对手那一次没发挥好。}加上脑电之后的最终成绩，三次都是 0.58 左右，没有任何变化。

三次平均之后，提升是 \key{-0.003}——也就是没有。

\key{我们最终没有采信这批结果。}但它们值得完整记录下来，理由有三个：

\begin{itemize}
\item 这是本项目中\key{客观测量最接近成功的一次}，如实记录比只报告结论更诚实；
\item 它清楚地展示了一个重要的方法学教训：\key{在样本不大的时候，看似漂亮的“提升”可能完全来自对手的波动}；
\item 如果未来要重新检验脑电，\key{应当从多种子平均开始}，而不是从单次训练的结果开始。
\end{itemize}

\section{三次“看起来成功了”：本阶段最重要的产出}\label{sec:false}

上一节的脑电种子问题不是孤例。这四天里，有三次我们的分析给出了阳性结果，如果直接采信，报告会写成“行为学、眼动、脑电深度模型分别发现了独立信号”。

\key{三次全是假的，三次全被拦下了。}

\begin{center}
\begin{tabularx}{\textwidth}{@{}lcY@{}}
\toprule
\textbf{阶段} & \textbf{涉及} & \textbf{真相} \\
\midrule
行为学 & 8 个结果 & 它们赢过的“年龄性别年级”基线，在那一小批孩子（342 人）里\key{本身低于抛硬币}。不是行为学变强了，是对手垮了。 \\
眼动 & 18 个结果 & 完全相同的机制，全部出现在同一台设备的小样本里。 \\
脑电深度 & 7 个结果 & 深度模型“赢过”传统特征——而传统特征在那个条件下也低于抛硬币。 \\
\bottomrule
\end{tabularx}
\end{center}

\subsection{为什么会这样}

打个比方：考试比的是两个学生的名次。如果对手因为生病考了零分，你赢了他，并不说明你考得好。

样本量小的时候，“年龄性别年级”这个对手会不稳定，有时会发挥失常到比瞎猜还差。这时任何一个平庸的模型都能“赢”它，而统计检验会一本正经地告诉你这个胜利是显著的。

\subsection{我们做了什么}

第一次（行为学）发现后，我们做了三项独立验证确认它是假的。最有说服力的一项是：\key{把孩子的诊断标签随机打乱 100 次重跑}——在完全没有真实关联的情况下，仍有约六分之一的次数能重现出同样大的“提升”。

然后给判断规则加了一条：\key{一个结果要被采信，它赢过的对手必须自己先站得住（高于抛硬币）。}

这条规则随后拦下了眼动的 18 个和脑电的 7 个。\key{如果没有它，这份报告会包含三个错误的阳性发现。}

\section{我们能确定什么，不能确定什么}

\subsection{可以确定的}

\begin{itemize}
\item 在\key{这一批孩子}身上、用\key{目前这套方法}，五种客观测量都不能超过年龄性别年级。
\item 这个结论不是数据质量或程序错误造成的——五种数据的生理效度都独立验证过。
\item 面部表情中确实存在超出“房间和长相”的真实信息，只是不足以超过人口学。
\item 眼动范式中的“情绪注意偏向”指标，在当前设计下测不准。
\end{itemize}

\subsection{不能由此推断的}

\begin{itemize}
\item \key{不能说“心理问题没有客观生物标记”。}我们检验的是特定的几种测量、特定的分析方法、特定的一批孩子，而且其中若干条已知存在采集或设计缺陷。
\item \key{不能说这些数据没有价值。}它们能很好地预测年龄和性别，在其他研究问题上仍然有用。
\item \key{不能说“再多试几种方法就一定能成”。}我们已经用了从传统特征到深度学习的多种方法。
\end{itemize}

\section{下一步的四个方向}\label{sec:next}

\begin{enumerate}
\item \key{反馈已发现的采集缺陷。}把赌博任务响应窗口设置、必可明按键未采集这两个问题反馈给设备方与实验设计方，避免在后续批次里重复发生。

\item \key{如果还有新一轮采集，修改实验设计。}情绪面孔从每种 12 张增加到 40 张以上，扫视任务从 8 个试次增加到 40 个以上，让每台设备在每所学校都采一部分。这三项改动直接针对第 \ref{sec:why} 节查出的原因。

\item \key{重新考虑研究问题的定义。}目前的“有／无心理问题”把抑郁、焦虑、高风险合并成一类，可能过于笼统。换成症状严重程度的连续评分，或某一类特定症状，客观测量可能有更好表现。建议与医院方共同讨论。

\item \key{把已验证的测量用于其他问题。}这些数据在年龄和性别上的预测能力很强（0.77--0.82）。发育轨迹、认知能力发展等方向可能比疾病分类更适合这批数据。
\end{enumerate}

\vspace{1em}
\noindent\rule{\textwidth}{0.4pt}
\vspace{0.4em}

\noindent\small 本报告为面向非技术读者的汇报版本，略去了方法细节、完整结果表和统计过程。逐项数据、方法说明与全部结果表见同目录下的技术版报告 \fname{chongqing\_progress\_20260907\_20260910\_technical.tex}（27 页），以及项目内的 \texttt{PROGRESS.md} 与 \texttt{reports/} 各阶段报告。本报告中的每一个数字均来自项目结果表，未作估算。

\end{document}
